\title{DeepSeekMath: Pushing the Limits of Mathematical Reasoning in Open Language Models}

\begin{abstract}
Language models encounter considerable difficulty with mathematical reasoning, primarily due to its intricate and highly structured characteristics. In this work, we present DeepSeekMath~7B, developed by further pre-training DeepSeek-Coder-Base-v1.5 7B using 120B mathematics-oriented tokens drawn from Common Crawl, in addition to related text and code. Without leveraging external toolkits or ensemble methods, DeepSeekMath~7B attains an outstanding 51.7\% on the competition-level MATH benchmark, closely rivaling state-of-the-art models such as Gemini-Ultra and GPT-4. Its self-consistency, evaluated across 64 sampled outputs, reaches 60.9\% on the same benchmark. The remarkable mathematical proficiency of DeepSeekMath~7B stems from two primary aspects: firstly, we fully utilize publicly accessible web resources via a rigorously designed data filtration framework; secondly, we propose Group Relative Policy Optimization (GRPO)—a novel adaptation of Proximal Policy Optimization (PPO)—which augments mathematical reasoning capabilities while also optimizing PPO’s memory efficiency.
\end{abstract}

\section{Introduction}

The advent of large language models (LLMs) has dramatically transformed how artificial intelligence systems address mathematical reasoning, leading to significant improvements in benchmarks focused on quantitative \citep{MATH} and geometric reasoning \citep{trinh2024solving}. Additionally, LLMs have emerged as valuable tools for supporting humans in tackling advanced mathematical challenges \citep{tao}. Nonetheless, despite the impressive results demonstrated by leading models such as GPT-4 \citep{gpt4} and Gemini-Ultra \citep{gemini}, these systems remain unavailable to the public. In contrast, the current generation of open-source alternatives falls notably short in terms of performance.

This paper presents DeepSeekMath, a specialized language model designed for mathematics that demonstrates substantial improvements over existing open-source models and closely rivals GPT-4 on academic evaluation sets. Central to this advancement is the construction of the DeepSeekMath Corpus, an expansive and meticulously curated pre-training dataset containing 120B mathematics-related tokens. The data are sourced from Common Crawl (CC) through a fastText-based classification system \citep{joulin2016fasttext}. For its initial deployment, the classifier leverages positive examples from OpenWebMath \citep{openwebmath} and utilizes a broad assortment of general web documents as negatives. Following this, the classifier is applied to the CC to uncover further positive examples, which are then refined and validated by human annotators. This revised and higher quality set of examples is used to retrain the classifier, leading to enhanced extraction performance. Evaluation demonstrates that this extensive corpus yields a superior model: DeepSeekMath-Base 7B achieves 64.2\% on GSM8K \citep{gsm8k} and 36.2\% on the challenging MATH benchmark \citep{MATH}, outperforming models such as Minerva 540B \citep{minerva}. Furthermore, as the DeepSeekMath Corpus encompasses multiple languages, noticeable gains are also observed on Chinese mathematical benchmarks \citep{wei2023cmath,agieval}. We regard our approach to constructing mathematical datasets as an initial step for further advancements in the research field, acknowledging ample opportunities for future progress.

DeepSeekMath-Base is constructed by initializing from DeepSeek-Coder-Base-v1.5 7B \citep{deepseek-coder}, as preliminary evidence suggests that using a code-oriented pre-trained model serves as a more robust foundation than a generic large language model (LLM). Additionally, we find that mathematical pre-training not only augments the model's proficiency in mathematical domains but also bolsters its performance on general benchmarks such as MMLU \citep{mmlu} and BBH \citep{bbh}, highlighting improvements in overall reasoning abilities beyond mathematics alone.

Subsequent to pre-training, we perform instruction tuning focused on mathematics for DeepSeekMath-Base. This stage incorporates data involving chain-of-thought \citep{cot}, program-of-thought \citep{pot,pal}, and tool-augmented reasoning \citep{tora}. The instruction-tuned model, DeepSeekMath-Instruct 7B, surpasses all other 7B-scale models and achieves performance on par with instruction-tuned models at the 70B scale within the open-source community.

Moreover, we propose Group Relative Policy Optimization (GRPO), a novel reinforcement learning (RL) method that modifies Proximal Policy Optimization (PPO) \citep{schulman2017proximal}. Unlike PPO, GRPO eliminates reliance on a separate critic network, deriving baselines instead from group-level scores, which yields significant savings in computational costs. Utilizing only a subset of English instruction-tuning data during RL training, GRPO delivers marked gains over the already strong DeepSeekMath-Instruct model. Improvements manifest across both in-domain tasks (GSM8K: 82.9\% $\rightarrow$ 88.2\%, MATH: 46.8\% $\rightarrow$ 51.7\%) and tasks outside the training distribution, such as CMATH (84.6\% $\rightarrow$ 88.8\%).

In addition, we offer a unified conceptual framework to interpret methods such as Rejection Sampling Fine-Tuning (RFT) \citep{yuan2023scaling}, Direct Preference Optimization (DPO) \citep{dpo}, PPO, and our proposed GRPO. This unified perspective categorizes these techniques as variants or simplified forms of RL. We conduct comprehensive analyses, comparing, for instance, online versus offline optimization, the effects of supervising the outcome as opposed to the process, and evaluating single-turn versus iterative RL approaches, in order to uncover fundamental components underlying this paradigm. Finally, we elucidate why reinforcement learning further enhances instruction-tuned models and summarize promising avenues for advancing RL efficacy under this unified framework.

\subsection{Contributions}
We contribute by advancing scalable math pre-training techniques and by conducting in-depth investigations into reinforcement learning.

\textbf{Math Pre-Training at Scale}

\begin{itemize}[topsep=0pt]
    \item Our study demonstrates that the publicly available Common Crawl dataset contains rich mathematical information. Through a carefully engineered data selection process, we compile the DeepSeekMath Corpus, a robust dataset containing 120B tokens extracted from web sources with a focus on mathematics. This corpus is nearly 7 times larger than the mathematical web content used by Minerva \citep{minerva} and exceeds the size of OpenWebMath \citep{openwebmath} by a factor of 9.

    \item The pre-trained DeepSeekMath-Base 7B model delivers results on par with Minerva 540B \citep{minerva}, illustrating that increasing parameter count alone does not guarantee superior mathematical reasoning. Our findings suggest that even models with fewer parameters can excel when pre-trained on highly curated datasets.

    \item We present insights from our math-focused training experiments. Pre-training on code before introducing math tasks enhances the model’s proficiency in mathematical problem-solving, both when employing tools and otherwise. This provides evidence toward resolving the longstanding query: \textit{does code training contribute to reasoning abilities?} Our results support this, particularly within the domain of mathematical reasoning.

    \item While the use of arXiv papers for model training is widespread, especially for mathematics-focused tasks, our experiments show that it yields no significant gains across the mathematical benchmarks considered in this work.
\end{itemize}

\textbf{Exploration and Analysis of Reinforcement Learning}

\begin{itemize}[topsep=0pt]
    \item We propose Group Relative Policy Optimization (GRPO), a reinforcement learning algorithm designed for both efficiency and efficacy. Unlike methods such as Proximal Policy Optimization (PPO), GRPO dispenses with the critic model, instead utilizing group scores to estimate the baseline. This shift enables a significant reduction in training demands.
    \item Through empirical study, we show that GRPO substantially improves the capabilities of our instruction-tuned model, DeepSeekMath-Instruct, relying solely on instruction-tuning datasets. Additionally, we find that applying GRPO yields notable gains in generalization to out-of-domain tasks during reinforcement learning.
    \item We articulate a coherent framework that connects several methodologies—RFT, DPO, PPO, and GRPO—under a unified perspective. Our investigation includes thorough experimental comparisons, such as online versus offline learning paradigms, outcome versus process-based supervision, and contrasts between single-turn and iterative reinforcement learning. These studies illuminate key drivers of the overarching approach.
    \item Drawing on insights from our unified framework, we analyze the fundamental factors underlying effective reinforcement learning, and outline future avenues that could lead to further advancements in training large language models with reinforcement learning techniques.
\end{itemize}

\subsection{Summary of Evaluations and Metrics}

\begin{itemize}[topsep=0pt]
    \item \textbf{English and Chinese Mathematical Reasoning}: Our models undergo comprehensive evaluation using a spectrum of English and Chinese benchmarks that encompass mathematics problems spanning elementary through collegiate difficulty. English test suites include GSM8K \citep{gsm8k}, MATH \citep{MATH}, SAT \citep{llemma}, OCW Courses \citep{minerva}, and MMLU-STEM \citep{mmlu}; for Chinese, we use MGSM-zh \citep{mgsm}, CMATH \citep{wei2023cmath}, Gaokao-MathCloze \citep{agieval}, and Gaokao-MathQA \citep{agieval}. Model evaluation is twofold: the ability to produce fully articulated written solutions unaided by tools, and performance in computational problem-solving using Python.

    On English datasets, DeepSeekMath-Base attains performance on par with the proprietary Minerva 540B model \citep{minerva} and consistently surpasses open-source base models such as Mistral 7B \citep{mistral} and Llemma-34B \citep{llemma}, irrespective of prior math-specific pre-training, often by substantial margins. Remarkably, DeepSeekMath-Base leads on Chinese tasks, which we attribute to our decision not to restrict pre-training exclusively to English mathematical corpora (as in previous studies \citep{minerva,llemma}), and instead to incorporate extensive, high-quality data from non-English sources. The addition of mathematical instruction tuning and reinforcement learning, resulting in DeepSeekMath-Instruct and DeepSeekMath-RL, yields robust model performance: notably, these models are the first from the open-source domain to exceed 50\% accuracy on the MATH competition-level benchmark.
\end{itemize}

\item \textbf{Formal Mathematics}: For evaluating DeepSeekMath-Base in formal mathematics, we utilize the informal-to-formal theorem proving task described by \citep{dsp_proof} on the miniF2F benchmark \citep{minif2f}, employing Isabelle \citep{isabelle} as the designated proof assistant. DeepSeekMath-Base exhibits notable performance in few-shot autoformalization.

\item \textbf{Natural Language Understanding, Reasoning, and Code}: To comprehensively assess the model's general proficiency in understanding, reasoning, and programming, DeepSeekMath-Base is benchmarked on Massive Multitask Language Understanding (MMLU) \citep{mmlu}, which features 57 multiple-choice questions from various fields. We also use BIG-Bench Hard (BBH) \citep{bbh}, composed of 23 tasks predominantly requiring multi-step logical reasoning, and HumanEval \citep{codex} and MBPP \citep{mbpp} for code generation evaluation. Results indicate that mathematical pre-training enhances both reasoning and general language understanding capabilities.

\section{Math Pre-Training}

\subsection{Data Collection and Decontamination} 
This section details our methodology for assembling the DeepSeekMath Corpus from the Common Crawl dataset. Figure \ref{fig:data_collect} outlines an iterative pipeline that illustrates the systematic expansion of a large-scale mathematics-specific corpus from Common Crawl, initiated by a carefully curated seed dataset. While focused on mathematics, this pipeline is equally adaptable to other specialized domains, including programming.

Our initial seed is OpenWebMath \citep{openwebmath}, which is composed of rigorously selected mathematical web documents. With this collection, we train a fastText classifier \citep{joulin2016fasttext} aimed at identifying additional web pages resembling OpenWebMath. For model training, we sample 500,000 items from the seed corpus as positives and another 500,000 randomly chosen web pages from Common Crawl as negatives. Training utilizes an open-source toolkit\footnote{\url{https://fasttext.cc}} with these settings: a vector dimension of 256, learning rate of 0.1, word n-grams up to a length of 3, a minimum word frequency threshold of 3, and 3 training epochs. To make Common Crawl more manageable, we apply both URL-based deduplication and near-duplicate removal, which reduces the pool to 40 billion HTML web pages. Subsequently, the fastText model retrieves mathematics-related web pages from this deduplicated subset. To maintain data quality, all candidate documents are ranked based on the fastText scores, and only those with the highest rankings are selected for retention. The effectiveness of different data volumes is empirically evaluated using pre-training trials on the top 40B, 80B, 120B, and 160B tokens; for the initial iteration, the top 40B tokens are retained.

Following the initial round of data gathering, a considerable number of mathematical web pages are left uncollected, primarily because the set of positive samples used to train the fastText model is insufficiently diverse. To address this limitation, we expand the seed dataset by incorporating more mathematical web sources, thereby enhancing the performance of the fastText classifier. In our procedure, we first partition the entirety of Common Crawl into separate domains, where each domain represents all web pages that share a common base URL. For every domain, we determine the proportion of web pages retrieved during the first collection cycle. Domains from which more than 10\% of pages have been obtained are tagged as math-oriented (for instance, \url{mathoverflow.net}).

Next, we manually annotate URLs pointing to mathematical content within these math-related domains (such as \url{mathoverflow.net/questions}). Any web pages linked to these annotated URLs that were not collected in earlier rounds are subsequently added to the seed dataset. This methodology supplies more positive samples for training, resulting in a refined fastText model capable of capturing a greater number of mathematical web pages in the subsequent iterations. After four iterations, this process yields 35.5 million mathematical pages, amounting to 120 billion tokens. By the fourth round, we observe that approximately 98\% of pages included had already been gathered in the third round, prompting us to terminate the data collection process at that point.

To guard against contamination of evaluation benchmarks, we apply the filtration strategy introduced by \cite{deepseek-coder}. This involves excluding any web pages containing questions or answers from notable English mathematical benchmarks such as GSM8K \citep{gsm8k} and MATH \citep{MATH}, as well as Chinese benchmarks including CMATH \citep{wei2023cmath} and AGIEval \citep{agieval}. Specifically, we remove from our math training data any textual fragment where a contiguous 10-gram substring exactly matches one found in these benchmarks. For benchmark entries shorter than 10 grams but at least 3 grams long, we conduct exact matching to excise any contaminated web pages.

\subsection{Validating the Quality of the DeepSeekMath~Corpus}

To assess the relative quality of the DeepSeekMath~Corpus, we conduct pre-training trials comparing it against several recently introduced mathematical corpora:

\begin{itemize}[topsep=0pt]
    \item \textbf{MathPile} \citep{mathpile}: a multi-source dataset (8.9B tokens) compiled from resources including textbooks, Wikipedia, ProofWiki, CommonCrawl, StackExchange, and arXiv—with over 85\% of the data originating from arXiv;
    \item \textbf{OpenWebMath} \citep{openwebmath}: a collection of mathematical content extracted from CommonCrawl, totaling 13.6B tokens;
    \item \textbf{Proof-Pile-2} \citep{llemma}: an aggregated mathematical dataset composed of OpenWebMath, AlgebraicStack (10.3B tokens of math code), and arXiv articles (28.0B tokens). For experiments utilizing Proof-Pile-2, we adopt the 2:4:1 ratio of arXiv:Web:Code prescribed by \cite{llemma}.
\end{itemize}

\subsubsection{Training Setting}
\label{sec:quality-policy}
We fine-tune a general-purpose language model comprising 1.3B parameters, which utilizes the same architecture as the DeepSeek LLMs \citep{deepseek-llm}, referred to here as DeepSeek-LLM 1.3B. Separate models are trained on individual mathematical corpora, each for a total of 150B tokens. All training procedures employ the HAI-LLM \citep{haillm} framework, recognized for its efficiency and lightweight implementation. Consistent with the methodology adopted for DeepSeek LLMs, the AdamW optimizer \citep{adamW} is employed, configured with $\beta_1=0.9$, $\beta_2=0.95$, and $\mathrm{weight\_decay}=0.1$. The learning rate follows a multi-phase schedule: it reaches its maximum after 2,000 warmup iterations, decays to 31.6\% of the peak value at the 80\% mark of training, and then drops further to 10.0\% by 90\% completion. The learning rate is capped at a maximum of 5.3e-4, with training performed using a batch size of 4M tokens and a context window of 4K tokens.

\subsubsection{Evaluation Results}

\textbf{The DeepSeekMath~Corpus distinguishes itself by its high quality, extensive multilingual content, and unprecedented scale.}

\begin{itemize}[topsep=0pt]
    \item \textbf{High-quality}: 
    We assess downstream task performance across eight mathematical benchmarks via few-shot chain-of-thought prompting \cite{cot}. Table \ref{tab:corpora_comparison} highlights the clear advantage of models trained with the DeepSeekMath~Corpus. As depicted in Figure \ref{fig:corpora_comparisons}, models trained on DeepSeekMath Corpus outperform those trained on Proof-Pile-2 at 50B tokens (equivalent to one complete pass through Proof-Pile-2), underscoring the superior average quality of DeepSeekMath Corpus.
    \item \textbf{Multilingual}:
    The DeepSeekMath~Corpus includes content spanning various languages, with English and Chinese making up the majority. As reported in Table \ref{tab:corpora_comparison}, leveraging the DeepSeekMath~Corpus during training results in enhanced mathematical reasoning capabilities for both English and Chinese inputs. Conversely, existing math corpora, which are predominantly English-only, provide minimal gains or may even be detrimental to Chinese mathematical reasoning proficiency.
    \item \textbf{Large-scale}:
    The DeepSeekMath~Corpus significantly exceeds the size of all previously available mathematical datasets. Figure \ref{fig:corpora_comparisons} illustrates that, with the DeepSeekMath~Corpus, DeepSeek-LLM 1.3B demonstrates more rapid and sustained improvement. In comparison, baseline corpora are far smaller in scale, have already undergone numerous training passes, and consequently show model performance leveling off quickly.
\end{itemize}

\subsection{Training and Evaluating DeepSeekMath-Base 7B}

This section details the development of DeepSeekMath-Base 7B, a foundational model engineered to excel in mathematical reasoning. Built upon DeepSeek-Coder-Base-v1.5 7B \citep{deepseek-coder}, our approach involves further training the model on a total of 500B tokens. The training data is distributed as follows: 56\% from the DeepSeekMath Corpus, 4\% from AlgebraicStack, 10\% from arXiv, 20\% consisting of code from Github, and the final 10\% comprised of natural language data sourced from Common Crawl in both English and Chinese. The majority of the training protocol aligns with the methodology described in Section \ref{sec:quality-policy}, but with adjustments in the maximum learning rate, which is set at 4.2e-4, and a batch size increased to 10M tokens.

A thorough evaluation is undertaken to gauge the model's mathematical proficiency. The assessment criteria include generating standalone mathematical solutions without auxiliary tools, utilizing computational tools for mathematical problem-solving, and engaging in formal theorem proving. To provide a holistic perspective, we also analyze the broader capabilities of the base model, including natural language understanding, reasoning, and programming.

\paragraph{Mathematical Problem Solving with Step-by-Step Reasoning}

We assess DeepSeekMath-Base's mathematical problem-solving aptitude using few-shot chain-of-thought prompting \citep{cot} on eight benchmarks in both English and Chinese. The selected benchmarks represent a range of quantitative reasoning tasks (such as GSM8K \citep{gsm8k}, MATH \citep{MATH}, and CMATH \citep{wei2023cmath}) as well as multiple-choice questions (such as MMLU-STEM \citep{mmlu} and Gaokao-MathQA \citep{agieval}), thus covering mathematical topics from elementary up to collegiate levels of difficulty.

As detailed in Table \ref{tab:base_math_cot}, among open-source base models, DeepSeekMath-Base 7B consistently achieves the highest scores across all eight benchmarks. This includes outperforming the widely adopted Mistral 7B \citep{mistral} and the newly launched Llemma 34B \citep{llemma}, which is specifically trained on the Proof-Pile-2 math corpus \citep{llemma}. Remarkably, on the MATH dataset, which assesses competition-level mathematical reasoning, DeepSeekMath-Base registers an absolute improvement exceeding 10\% compared to any existing open-source base model. Furthermore, it exceeds the performance of Minerva 540B \citep{minerva}, a proprietary base model that is 77 times larger, developed on PaLM \citep{palm} and further specialized on mathematical content.

\paragraph{Mathematical Problem Solving with Tool Use} To assess program-supported mathematical reasoning, we employ few-shot program-of-thought prompting \citep{pot,pal} on both GSM8K and MATH. Here, models are prompted to tackle each question by producing a Python script, utilizing libraries like \textit{math} and \textit{sympy} for sophisticated calculations. The answer is determined by executing the script. Table \ref{tab:base_math_pot_proof} indicates that DeepSeekMath-Base 7B surpasses the former leading Llemma 34B model.

\paragraph{Formal Mathematics}

Automating formal proofs plays an important role in verifying the precision and dependability of mathematical reasoning, thereby increasing efficiency and attracting growing research focus. We measure the abilities of DeepSeekMath-Base 7B on the informal-to-formal proof generation task from \citep{dsp_proof}, which requires generating a formal proof using an informal proposition, its formal translation, and an accompanying informal argument. This evaluation is carried out on miniF2F \citep{minif2f}, a benchmark tailored for Olympiad-grade formal mathematics, prompting the model to craft an Isabelle formalization for each question in a few-shot setting. Consistent with the methodology of \cite{dsp_proof}, the model produces proof outlines which are then completed by Sledgehammer, an existing automated theorem prover \citep{sledgehammer}. Table \ref{tab:base_math_pot_proof} demonstrates that DeepSeekMath-Base 7B attains strong outcomes for automated proof formalization.

\paragraph{Natural Language Understanding, Reasoning, and Code}

To gauge capabilities in natural language processing, reasoning, and programming, we test DeepSeekMath-Base 7B on MMLU \citep{mmlu} for language understanding, BBH \citep{bbh} for reasoning, and HumanEval \citep{codex} along with MBPP \citep{mbpp} for code generation. Results in Table \ref{tab:understanding_reasoning_code} show marked gains over its predecessor, DeepSeek-Coder-Base-v1.5 \citep{deepseek-coder}, on MMLU and BBH, highlighting the beneficial effects of specialized mathematical training. Moreover, by continuing to train with coding data, DeepSeekMath-Base 7B sustains the previous model's performance across the coding benchmarks. Across the three reasoning and programming tasks, DeepSeekMath-Base 7B also substantially outperforms the generic Mistral 7B \citep{mistral}.

\section{Supervised Fine-Tuning}

\subsection{SFT Data Curation}

To develop our mathematical instruction-tuning dataset, we gathered both English and Chinese math problems representing a range of domains and difficulty levels. Each problem is paired with detailed solutions utilizing chain-of-thought (CoT) \citep{cot}, program-of-thought (PoT) \citep{pot,pal}, and tool-integrated reasoning formats \citep{tora}. The curated dataset comprises 776K training samples in total.

\begin{itemize}[topsep=0pt]
    \item \textbf{English mathematical datasets}: We enriched GSM8K and MATH with tool-integrated solutions and incorporated select examples from MathInstruct \citep{MathInstruct} and the Lila-OOD \citep{lila} training set, both employing CoT and PoT approaches. Our English corpus spans numerous mathematical branches, including but not limited to algebra, probability, number theory, calculus, and geometry.
    \item \textbf{Chinese mathematical datasets}: We amassed a dataset of K-12 math questions in Chinese, encompassing 76 specific subtopics such as linear equations, with answer annotations provided in CoT as well as tool-integrated reasoning formats.
\end{itemize}

\subsection{Training and Evaluating DeepSeekMath-Instruct 7B}

Here, we detail the instruction-tuning of DeepSeekMath-Instruct 7B, starting from DeepSeekMath-Base as the foundational model. During supervised fine-tuning, we randomly concatenate training instances, maintaining a maximum context length of 4K tokens. The training proceeds for 500 steps with a batch size set to 256, using a fixed learning rate of 5e-5.

To assess the mathematical problem-solving capabilities of our models, both with and without external tool support, we employ four quantitative reasoning benchmarks in both English and Chinese. We further position our results alongside current state-of-the-art models:

\begin{itemize}[topsep=0pt]
    \item \textbf{Closed-source models}: Benchmarks include several notable proprietary models: (1) The GPT series, notably GPT-4 \citep{gpt4} and GPT-4 Code Interpreter~\footnote{\url{https://openai.com/blog/chatgpt-plugins\#\#code-interpreter}}, (2) Gemini Ultra and Pro \citep{gemini}, (3) Inflection-2 \citep{inflection-2}, (4) Grok-1~\footnote{\url{https://x.ai/model-card}}, and prominent releases from Chinese tech firms including (5) Baichuan-3~\footnote{\url{https://www.baichuan-ai.com}}, as well as (6) the latest GLM-4~\footnote{\url{https://open.bigmodel.cn/dev/api\#glm-4}}

    from the GLM family \citep{glm}.
\end{itemize}

The majority of these models are designed for broad applicability and have typically undergone multiple rounds of alignment.

\item \textbf{Open-source models} comprise: (1) DeepSeek-LLM-Chat 67B \citep{deepseek-llm}, (2) Qwen 72B \citep{qwen}, (3) SeaLLM-v2 7B \citep{seallm}, and (4) ChatGLM3 6B \citep{chatglm3}, which serve as general-purpose models, as well as a set of models featuring mathematical capability enhancements, including (5) InternLM2-Math 20B~\footnote{\url{https://github.com/InternLM/InternLM-Math}}, which augments InternLM2 with specialized math instruction tuning, (6) Math-Shepherd-Mistral 7B, which leverages PPO training \citep{schulman2017proximal} on the Mistral 7B \citep{mistral} architecture utilizing a reward model based on supervised processes, (7) WizardMath series \citep{wizardmath}, which bolsters mathematical reasoning in Mistral 7B and Llama-2 70B \citep{llama2} through evolve-instruct methodology (a variant of instruction tuning using instructions evolved by AI) and PPO training on datasets predominantly drawn from GSM8K and MATH, (8) MetaMath 70B \citep{metamath}, which involves fine-tuning Llama-2 70B on an enriched compilation of GSM8K and MATH, (9) ToRA 34B \cite{tora}, which adapts CodeLlama 34B for tool-assisted mathematical reasoning through targeted fine-tuning, and (10) MAmmoTH 70B \citep{MathInstruct}, which comprises Llama-2 70B further trained via MathInstruct's instructional approach.
\end{itemize}

Table \ref{tab:sft_rl_math} presents results indicating that, when evaluated without tool integration, DeepSeekMath-Instruct 7B achieves impressive proficiency in multi-step reasoning tasks. On the challenging MATH dataset, the model consistently exceeds the performance of all current open-source models as well as the vast majority of proprietary models (including Inflection-2 and Gemini Pro), achieving at least a 9\% absolute margin. This advantage persists even in comparison to much larger models, such as Qwen 72B, or those enhanced by math-specific reinforcement learning techniques like WizardMath-v1.1 7B. While DeepSeekMath-Instruct is competitive with proprietary Chinese models GLM-4 and Baichuan-3 on the MATH benchmark, it still lags behind GPT-4 and Gemini Ultra.

When the evaluation protocol permits models to combine natural language-based reasoning with programmatic tool use for problem-solving, DeepSeekMath-Instruct 7B approaches a 60\% accuracy rate on MATH, outperforming every other open-source model assessed. Furthermore, on other evaluation benchmarks, its performance rivals that of DeepSeek-LLM-Chat 67B, a previous state-of-the-art model whose parameter count is tenfold higher.

\section{Reinforcement Learning}

\subsection{Group Relative Policy Optimization} Reinforcement learning (RL) has consistently demonstrated its utility in elevating the mathematical reasoning skills of large language models after the Supervised Fine-Tuning (SFT) phase \citep{wang2023math,wizardmath}. In this context, we present Group Relative Policy Optimization (GRPO), an RL approach characterized by its efficiency and effectiveness.

\subsubsection{From PPO to GRPO} Proximal Policy Optimization (PPO) \citep{schulman2017proximal} remains the predominant actor-critic RL technique employed during the fine-tuning phase of LLMs \citep{ouyang2022training}. PPO seeks to maximize a surrogate objective defined as follows:

\begin{equation}
\footnotesize
    \mathcal{J}_{PPO}(\theta) = \mathbb{E}{[q \sim P(Q), o \sim \pi_{\theta_{old}}(O|q)]} \frac{1}{|o|} \sum_{t=1}^{|o|} \min \left[ \frac{\pi_\theta(o_{t} | q, o_{<t})}{\pi_{\theta_{old}}(o_{t} | q, o_{<t})} A_{t}, \text{clip} \left( \frac{\pi_\theta(o_{t} | q, o_{<t})}{\pi_{\theta_{old}}(o_{t} | q, o_{<t})

Here, $\pi_{\theta}$ denotes the current policy, while $\pi_{\theta_{old}}$ represents the previous iteration of the policy model. The variables $q$ and $o$ correspond to questions and generated outputs, which are sampled from a dataset of questions and from the earlier policy $\pi_{\theta_{old}}$, respectively. The hyper-parameter $\varepsilon$ is introduced in PPO as a means to stabilize the training process by clipping. The term $A_t$ refers to the advantage function, which is calculated using Generalized Advantage Estimation (GAE) \citep{gae}, utilizing the set of rewards $\{r_{\ge t}\}$ and a learned value estimator $V_{\psi}$. Consequently, PPO requires training a separate value network in tandem with the policy model. To prevent the reward model from being excessively exploited, it is standard to include a KL-divergence-based penalty per token in the reward, drawing from a reference model as outlined in \citep{ouyang2022training}. This reward formulation is given by:

\begin{equation}
    r_{t} = r_\varphi(q, o_{\le t}) - \beta \log\frac{\pi_{\theta}(o_{t}|q, o_{<t})}{\pi_{ref}(o_{t}|q, o_{<t})},
\label{eq:PPO-reward}
\end{equation}

where $r_\varphi$ represents the reward model, $\pi_{ref}$ is the reference model (commonly set as the initial SFT model), and $\beta$ is the scaling parameter for the KL penalty.

Training a value network in PPO, which generally matches the policy model in scale, imposes considerable memory and computational demands. Furthermore, in reinforcement learning, the value network acts as a baseline in advantage estimation to help reduce variance. However, in LLM scenarios, the reward model often assigns a reward only to the final token, creating challenges for learning a well-calibrated per-token value function. To circumvent these issues, as illustrated in Figure \ref{fig:grpo}, we introduce Group Relative Policy Optimization (GRPO). This method eliminates the requirement for explicit value function modeling, and instead, leverages the mean reward across a collection of sampled outputs for the same question as a baseline. Specifically, for each question $q$, GRPO generates a group of outputs $\{o_1, o_2, \cdots, o_G\}$ sampled from the prior policy $\pi_{\theta_{old}}$, and then maximizes the following objective to train the policy:

\begin{equation}
\footnotesize
\begin{split}
    \mathcal{J}_{GRPO}(\theta) &= \mathbb{E}{[q \sim P(Q), \{o_i\}_{i=1}^G \sim \pi_{\theta_{old}}(O|q)]}  \\
    & \frac{1}{G}\sum_{i=1}^G\frac{1}{|o_i|} \sum_{t=1}^{|o_i|} \left\{ \min \left[ \frac{\pi_\theta(o_{i,t} | q, o_{i,<t})}{\pi_{\theta_{old}}(o_{i,t} | q, o_{i,<t})} \hat{A}_{i,t}, \text{clip} \left( \frac{\pi_\theta(o_{i,t} | q, o_{i,<t})}{\pi_{\theta_{old}}(o_{i,t} | q, o_{i,<t})}, 1 - \varepsilon, 1 + \varepsilon \right)  \hat{A}_{i,t} \right] - \beta \mathbb{D}_{KL}\left[\pi_{\theta} || \pi_{ref}\right]\right\} ,
\end{split}
\label{eq:GRPO-obj}
\end{equation}

with $\varepsilon$ and $\beta$ as the relevant hyper-parameters. Here, $\hat{A}_{i,t}$ refers to the advantage computed solely based on the relative rewards within each sampled group, the specifics of which are explained in the subsequent subsections. This group-relative formulation in GRPO resonates with the inherently comparative design of most reward models, which are commonly trained using paired or ranked outputs on shared prompts. It is also important to highlight that, contrary to the approach in PPO which adds the KL term directly to the reward, GRPO applies a KL-divergence regularization between the learned and reference policies in the loss term, thereby streamlining the computation of $\hat{A}_{i,t}$. Unlike the KL penalty formulation used in (\ref{eq:PPO-reward

\begin{algorithm}[t]
  \small
  \caption{Iterative Group Relative Policy Optimization}
  \textbf{Input} initial policy model $\pi_{\theta_{\text{init}}}$; reward models $r_\varphi$; task prompts $\mathcal{D}$; 
  hyperparameters $\varepsilon$, $\beta$, $\mu$
  \begin{algorithmic}[1]
    \State Set current policy $\pi_\theta \leftarrow \pi_{\theta_{\text{init}}}$
    \For{iteration = 1, \dots, I}
       \State Store reference model $\pi_{ref} \leftarrow \pi_{\theta}$
      \For{step = 1, \dots, M}
      \State Randomly select batch $\mathcal{D}_b$ from $\mathcal{D}$
      \State Assign old policy $\pi_{\theta_{old}} \leftarrow \pi_{\theta}$ 
      \State For each prompt $q$ in $\mathcal{D}_b$, generate $G$ responses $\{o_i\}_{i=1}^G$ sampled from $\pi_{\theta_{old}} (\cdot \mid q) $
      \State Calculate rewards $\{r_i\}_{i=1}^{G}$ for each $o_i$ by evaluating $r_{\varphi}$ 
      \State Use group-based estimation to obtain $\hat{A}_{i,t}$ for token $t$ of $o_i$
      \For{GRPO iteration = 1, \dots, $\mu$}
        \State Update policy $\pi_{\theta}$ by optimizing the GRPO criterion (Equation \ref{eq:GC-GRPO})
      \EndFor
    \EndFor \State Adjust $r_\varphi$ using continued training and replay data. 
    \EndFor 
  \end{algorithmic}
  \textbf{Output} $\pi_\theta$
  \label{alg:iter-grpo}
\end{algorithm}

\subsubsection{Outcome Supervision RL with GRPO} More formally, for a given question $q$, a collection of responses $\{o_1, o_2, \cdots, o_G\}$ is drawn from the prior policy $\pi_{\theta_{old}}$. The reward model then assigns a score to each response, producing a set of $G$ reward values $\mathbf{r}=\{r_1, r_2, \cdots, r_G\}$. These reward values are standardized by subtracting their mean and dividing by their standard deviation within the group. In the outcome supervision scenario, the standardized reward is allocated to the terminal token of each response $o_i$, and all tokens in $o_i$ receive this same advantage signal, specifically, $\hat{A}_{i, t} = \widetilde{r}_i = \frac{r_i- {\rm mean}(\mathbf{r})}{{\rm std}(\mathbf{r})}$. The policy is then trained to maximize the objective specified in equation (\ref{eq:GRPO-obj}).

\subsubsection{Process Supervision RL with GRPO} While outcome supervision offers feedback solely at the conclusion of an output, this granularity may not be adequate for guiding policies on intricate mathematical problems. As proposed in \cite{wang2023math}, we also consider process supervision, which supplies a reward after each step of reasoning. Concretely, given a question $q$ and $G$ generated solutions $\{o_1, o_2, \cdots, o_G\}$, a stepwise reward model is employed to score each intermediate reasoning segment, resulting in reward sequences: $\mathbf{R} = \{ \{r_1^{index(1)},\cdots,r_1^{index(K_1)}\}, \cdots,  \{r_G^{index(1)},\cdots,r_G^{index(K_G)}\} \}$, where $index(j)$ identifies the ending token position of the $j$-th step, and $K_i$ specifies the total number of steps for output $i$. These stepwise rewards are normalized within the group as $\widetilde{r}_i^{index(j)} = \frac{r_i^{index(j)} - {\rm mean}(\mathbf{R})}{{\rm std}(\mathbf{R})}}$. In this paradigm, the advantage for each token corresponds to the sum of normalized rewards from all subsequent steps: $\hat{A}_{i, t} = \sum_{index(j) \ge t} \widetilde{r}_i^{index(j)}$. The policy optimization again seeks to maximize the objective in equation (\ref{eq:GRPO

\subsubsection{Iterative RL with GRPO} As reinforcement learning advances, the previously trained reward model might no longer provide adequate supervision for the evolving policy model. To address this, we investigate the use of iterative RL utilizing GRPO. As outlined in Algorithm \ref{alg:iter-grpo}, iterative GRPO involves constructing fresh training data for the reward model by leveraging the samples produced by the current policy model. The reward model is then updated iteratively via a replay mechanism, which retains 10\% of the historical dataset. Subsequently, the policy model is adopted as the reference model, and further training is performed on the policy model using the updated reward model.

\subsection{Training and Evaluating DeepSeekMath-RL}

Our RL experiments are conducted on DeepSeekMath-Instruct 7B. For RL training, we use chain-of-thought-style questions that overlap with GSM8K and MATH present in the SFT dataset, totaling approximately 144,000 examples. We deliberately omit other SFT-based questions to clearly analyze the effect of RL on benchmarks that have not been directly exposed to data during the RL process. The reward model training set is built as described in \citep{wang2023math}. The initial reward model is finetuned from DeepSeekMath-Base 7B, employing a learning rate of 2e-5. In the GRPO procedure, the policy model is optimized with a learning rate of 1e-6 and a KL penalty coefficient set to 0.04. For each prompt, $64$ completions are generated. We fix the sequence length limit at 1024 tokens, with a batch size of 1024 during training. The policy model undergoes a single parameter update after each exploration round. We assess the performance of DeepSeekMath-RL 7B using the same benchmarks as DeepSeekMath-Instruct 7B. For evaluation purposes, tasks from GSM8K and MATH using chain-of-thought approaches are classified as in-domain, while all remaining benchmarks are treated as out-of-domain tasks.

Table \ref{tab:sft_rl_math} presents a comparative analysis of both open- and closed-source models employing chain-of-thought and tool-augmented reasoning techniques on English and Chinese evaluation sets. Our observations indicate that: (1) DeepSeekMath-RL 7B achieves accuracy rates of 88.2\% on GSM8K and 51.7\% on MATH when using chain-of-thought approaches, outperforming every open-source model ranging from 7B to 70B parameters, as well as surpassing most closed-source systems. (2) Notably, the DeepSeekMath-RL 7B model is exclusively fine-tuned on chain-of-thought-formatted instruction data derived from GSM8K and MATH, with DeepSeekMath-Instruct 7B serving as the initialization. Even with this narrow training regime, it consistently exceeds the performance of DeepSeekMath-Instruct 7B across all measured metrics, providing strong evidence for the effectiveness of reinforcement learning in this context.

\section{Discussion}
This section details insights drawn from our pre-training and reinforcement learning (RL) experimental processes.

\subsection{Lessons Learnt in Pre-Training}

We begin by summarizing our findings from the pre-training phase. Unless mentioned otherwise, all experiments adopt the configurations delineated in Section \ref{sec:quality-policy}. Importantly, any mention of the DeepSeekMath Corpus in this part refers specifically to the 89B-token dataset assembled during the project's second data collection cycle.

\subsubsection{Code Training Benefits Mathematical Reasoning}

It has long been hypothesized—though seldom rigorously tested—that training language models on code data enhances their reasoning skills. Our investigations provide empirical support for this claim within the realm of mathematical reasoning: incorporating code during training substantially improves models' mathematical problem-solving abilities, whether or not external tools are involved.

To rigorously assess the influence of code-based training on mathematical reasoning proficiency, we implemented and evaluated both a two-stage and a one-stage training paradigm:

\textbf{Two-Stage Training}

\begin{itemize}[topsep=0pt]
    \item \textbf{Code Training for 400B Tokens $\rightarrow$ Math Training for 150B Tokens}: DeepSeek-LLM 1.3B undergoes pre-training on 400B code tokens, after which it is further trained with 150B math tokens;
    \item \textbf{General Training for 400B Tokens $\rightarrow$ Math Training for 150B Tokens}: As a baseline, we replace the initial code data with 400B general tokens sampled from DeepSeek-AI’s large-scale corpus, aiming to analyze whether pre-training on code offers unique advantages for mathematical reasoning compared to generic textual data.
\end{itemize}

\textbf{One-Stage Training}

\begin{itemize}[topsep=0pt]
    \item \textbf{Math Training for 150B Tokens}: The model, DeepSeek-LLM 1.3B, is trained on 150B tokens from math datasets alone;
    \item \textbf{Training on a mixture of 400B Code Tokens and 150B Math Tokens}: Since sequential math training after code pre-training can harm coding capabilities, we also test if integrating code and math tokens during a single-stage training phase can not only foster mathematical reasoning skills but also prevent substantial degradation in coding abilities.
\end{itemize}

\paragraph{Results}
Performance metrics from Table \ref{tab:code-math} and Table \ref{tab:code-math-general-eval} provide a comprehensive comparison of each training protocol.

Incorporating code token training enhances program-based mathematical problem solving, applicable in both sequential and unified training regimes. As highlighted in Table \ref{tab:code-math}, code token pre-training alone yields substantial improvements on tasks such as GSM8K and MATH when solutions are derived using Python, with further benefits observed after subsequent math-focused fine-tuning. Additionally, simultaneous exposure to both code and math tokens during one-stage training substantially reduces catastrophic forgetting typical in two-stage settings, while also reinforcing both coding skills (Table \ref{tab:code-math-general-eval}) and program-assisted mathematical problem solving (Table \ref{tab:code-math}).

Furthermore, exposure to code tokens proves beneficial for mathematical reasoning even in scenarios without external tools. In the two-stage configuration, code pre-training confers noticeable advantages and accelerates the positive effects of math-focused training, culminating in the strongest observed outcomes. However, directly combining code and math data for one-stage training appears to negatively affect mathematical reasoning when tool assistance is not leveraged. We hypothesize this may be due to capacity limitations of DeepSeek-LLM 1.3B, which might restrict the model’s ability to internalize both types of knowledge in a single training phase.

\subsubsection{ArXiv Papers Appear to Provide Limited Benefit for Mathematical Reasoning}

It is standard practice to incorporate arXiv papers into math pre-training datasets \citep{minerva,gpt-f,llemma,mathpile}, yet the specific contribution of this data source to mathematical reasoning skills has not been rigorously assessed. Somewhat surprisingly, our findings suggest that arXiv papers do not substantially enhance mathematical reasoning. To investigate this, we utilized a range of model sizes, such as DeepSeek-LLM 1.3B and DeepSeek-Coder-Base-v1.5 7B \citep{deepseek-coder}, and applied various arXiv datasets subjected to distinct preprocessing steps:

\begin{itemize}[topsep=0pt]
    \item \textbf{MathPile} \citep{mathpile}: an 8.9-billion-token collection, mostly scientific arXiv papers (over 85\%), assembled through cleaning and filtering heuristics;
    \item \textbf{ArXiv-RedPajama} \citep{redpajama}: the complete set of arXiv LaTeX source files with extraneous material (preambles, macros, comments, bibliographies) excised, yielding 28.0 billion tokens.
\end{itemize}

Our experiments involved training DeepSeek-LLM 1.3B on 150B tokens and DeepSeek-Coder-Base-v1.5 7B on 40B tokens from each arXiv-based dataset. The results consistently indicate that exclusive training on arXiv data does not foster improvement in mathematical reasoning: in fact, both models failed to show gains and sometimes exhibited degraded performance on a variety of math benchmarks spanning different levels of challenge. These included quantitative problem sets like GSM8K and MATH (see Table \ref{tab:arxiv-cot}), multiple-choice assessments such as MMLU-STEM (see Table \ref{tab:arxiv-cot}), as well as formal mathematics tasks evaluated by miniF2F (Table \ref{tab:arxiv_atp}).

Nevertheless, our conclusions should be regarded as provisional. Several aspects remain unexamined:

\begin{itemize}[topsep=0pt]
    \item Possible utility of arXiv data for specialized mathematical problems outside the current evaluation scope, such as translating formal proofs to their informal narrative forms;
    \item The influence of integrating arXiv data with other corpora;
    \item The prospect that larger-scale models may realize benefits from arXiv material not observed at smaller scales.
\end{itemize}

Addressing these questions will require further research, which we defer to future investigations.

\subsection{Insights of Reinforcement Learning}
\subsubsection{Toward a Unified Framework} Here, we introduce a general framework to analyze various training methodologies, including SFT, RFT, DPO, PPO, and GRPO, and experimentally examine factors shaping this unified view. Broadly speaking, the gradient of a given training technique with respect to model parameters $\theta$ can be formulated as follows:

\begin{equation}
    \nabla_{\theta}\mathcal{J}_{\textcolor{red}{\mathcal{A}}}(\theta) = \mathbb{E}[\underbrace{(q,o) \sim \textcolor{red}{\mathcal{D}}}_{Data \ Source}]\left( \frac{1}{|o|} \sum_{t=1}^{|o|}  \underbrace{GC_{{\mathcal{A}}}(q, o, t, \textcolor{red}{\pi_{{rf}}})}_{Gradient \ Coefficient}  \nabla_{\theta}\log \pi_{\theta}(o_t | q, o_{<t})\right).
\label{eq:objective}
\end{equation}

This formulation involves three principal elements: 1) \textit{Data Source $\mathcal{D}$}, specifying the dataset utilized for training; 2) \textit{Reward Function $\pi_{{rf}}$}, providing the feedback or reward signals used in training; and 3) \textit{Algorithm $\mathcal{A}$}, which transforms the training data and reward signals into a gradient coefficient $GC$ responsible for scaling penalties or reinforcements. The following widely used approaches are considered within this unified framework:

\begin{itemize}
    \item  \textbf{Supervised Fine-tuning (SFT)}: In SFT, a pretrained model undergoes additional training on SFT data that is curated by humans.
    \item \textbf{Rejection Sampling Fine-tuning (RFT)}: RFT adapts the SFT model further using filtered outputs that are generated from the SFT model in response to SFT questions; the filtration process retains only those outputs deemed correct.
    \item \textbf{Direct Preference Optimization (DPO)}: DPO refines the SFT model by training it with enhanced outputs sampled from itself, where a pairwise DPO loss function is applied.
    \item \textbf{Online Rejection Sampling Fine-tuning (Online RFT)}: In contrast to standard RFT, Online RFT initializes the policy with the SFT model but then continues to improve it using augmented outputs that are generated dynamically from the up-to-date policy.
    \item \textbf{PPO/GRPO}: PPO/GRPO begin with the SFT model as the initial policy and strengthen it further with outputs sampled from the evolving policy in real time.
\end{itemize}

Table \ref{tab:data_policy} presents a summary of the elements comprising these methods. For a comprehensive account of the derivation procedure, consult Appendix \ref{app:analysis-rl}.

\paragraph{Observation about Data Source} We categorize the origin of data as either online sampling or offline sampling. Online sampling refers to situations where training samples are obtained from the ongoing exploratory outputs of the current policy model, whereas offline sampling involves drawing data from the initial SFT model’s sampled results. Among the evaluated algorithms, RFT and DPO operate under the offline regime, while Online RFT and GRPO utilize online data collection.

Analysis of Figure \ref{fig:rl-analysis} reveals that Online RFT yields marked improvements over RFT on both benchmarks. During the initial phase of training, Online RFT's performance aligns closely with RFT, but as training proceeds, Online RFT establishes a clear lead, highlighting the strengths of online sampling. This finding aligns with expectations, as the early stages show considerable similarity between the actor and SFT model, causing their sampled data to differ only slightly. In subsequent training phases, however, the divergence in actor-sampled data becomes pronounced, making real-time data collection increasingly advantageous.

\paragraph{Observation about Gradient Coefficient} To adjust the model parameters, the algorithm transforms the reward signal into a gradient coefficient. In this work, we distinguish between reward functions derived from `Rule' and from `Model'. The `Rule' approach assesses response quality directly according to answer accuracy, whereas the `Model' approach involves a reward model trained to evaluate responses, itself built upon rule-based judgments. Equations \ref{eq:GC-RFT} and \ref{eq:GC-GRPO} underscore an essential distinction between GRPO and Online RFT: GRPO modifies its gradient coefficient in response to the magnitude of rewards assigned by the reward model, thereby enabling more nuanced reinforcement or penalization depending on the response. Conversely, Online RFT does not feature such reward magnitude-based adjustment; it simply rewards all correct responses equally, without penalizing those that are incorrect.

As illustrated in Figure \ref{fig:rl-analysis}, GRPO outperforms online RFT, thereby emphasizing the effectiveness of modifying positive and negative gradient coefficients. Moreover, GRPO+PS yields better outcomes than GRPO+OS, demonstrating the advantages of employing fine-grained, step-sensitive gradient coefficients. We also investigate iterative RL by performing two cycles of iteration in our experiments. The results, displayed in Figure \ref{fig:iter-rl}, reveal that iterative RL markedly enhances performance, with the most substantial gains observed during the initial iteration.

\subsubsection{Why RL Works?} Our approach applies reinforcement learning to a selected subset of instruction tuning data, resulting in considerable improvements over the baseline instruction-tuned model. To elucidate the mechanism behind these improvements, we assess both Pass@K and Maj@K accuracy of the Instruct and RL-enhanced models across two evaluation datasets. Figure \ref{fig:maj-pass} indicates that RL leads to significant gains in Maj@K, whereas Pass@K remains largely unaffected. These results suggest that reinforcement learning strengthens the robustness of the output distribution, primarily by increasing the likelihood of a correct answer within the TopK responses rather than by enhancing the model’s fundamental abilities. This observation aligns with findings from \citep{wang2023making}, which describe a \textbf{misalignment problem} for reasoning tasks in SFT models and demonstrate that applying preference alignment strategies \citep{yuan2023rrhf,song2023preference,wang2023making} can improve their reasoning abilities.

\subsubsection{How to Achieve More Effective RL?}
\label{sec:effec_rl}
We show that reinforcement learning is highly effective for mathematical reasoning tasks. Additionally, we propose a unified framework for understanding a range of prevalent training approaches, framing them all as variations of direct or simplified RL. According to Equation \ref{eq:objective}, three foundational components emerge in this framework: Data Source, Algorithm, and Reward Function. We suggest several prospective research directions for each of these components.

\paragraph{Data Source} The data source forms the foundation for all training approaches. In the context of RL, this refers specifically to unlabeled questions paired with responses sampled from the policy model. Our experiments are limited to prompts from the instruction tuning phase and utilize basic nucleus sampling for generating outputs. We hypothesize that this may account for why our RL setup predominantly boosts Maj@K. Future work will investigate applying our RL procedure to prompts outside the original training distribution and will incorporate \textbf{advanced sampling (decoding) strategies} such as tree-search-based approaches \citep{yao2023tree}. Furthermore, the adoption of \textbf{efficient inference techniques} \citep{xia-etal-2023-speculative,leviathan2023fast,kwon2023efficient,xia2024unlocking}, which directly impact the exploration capacity of policy models, is anticipated to be vital.

\paragraph{Algorithms} The algorithms leverage both the dataset and the feedback from the reward signal to calculate gradient coefficients, thereby updating the parameters of the model. Drawing from Equation \ref{eq:objective}, it is apparent that current approaches universally and fundamentally \textbf{TRUST} the output of the reward function to modulate the conditional likelihood of specific tokens, either amplifying or suppressing it. Nevertheless, it is unrealistic to expect the reward signal to be unfailingly accurate, particularly when dealing with highly sophisticated tasks. For instance, the PRM800K datasets \citep{lightman2023let}, though meticulously labeled by expert annotators, still present around 20\% erroneous annotations\footnote{\url{https://github.com/openai/prm800k/issues/12\#issuecomment-1728491852}}. This motivates the investigation of reinforcement learning algorithms with built-in robustness to unreliable or noisy reward signals. We anticipate that such \textbf{WEAK-TO-STRONG} \citep{burns2023weak} alignment frameworks may significantly transform the core mechanics of learning algorithms.

\paragraph{Reward Function} The reward function delivers the essential training signals throughout the RL process, commonly implemented via a neural reward model. We identify three pivotal avenues for advancing reward models: 1) \textbf{Improving the generalization of reward models.} The capacity to generalize to unseen queries and sophisticated decoding scenarios is critical, as inadequate generalization may lead RL processes to simply reinforce existing LLM distributions instead of cultivating deeper abilities; 2) \textbf{Incorporating reward model uncertainty.} Explicitly modeling uncertainty offers a mechanism for integrating weak reward signals with weak-to-strong learning paradigms; 3) \textbf{Developing scalable and high-fidelity process reward models} that can yield detailed, process-level feedback for reasoning-oriented tasks \citep{lightman2023let,wang2023math}.

\section{Conclusion, Limitation, and Future Work}

In this work, we introduce DeepSeekMath, which sets a new standard among open-source models on the competition-grade MATH benchmark, coming close to matching the capabilities of proprietary systems. DeepSeekMath~builds upon DeepSeek-Coder-v1.5 7B, and is further continuously trained over 500B tokens, notably leveraging a large subset—120B math-specific tokens—drawn from Common Crawl. Through a thorough ablation study, we demonstrate that web pages are a particularly promising resource for high-quality mathematical content, whereas the contributions from arXiv are less significant than anticipated. We also propose Group Relative Policy Optimization (GRPO), an adaptation of Proximal Policy Optimization (PPO), which allows for improved mathematical reasoning proficiency while reducing memory usage. Our experimental evaluations confirm the efficacy of GRPO even for DeepSeekMath-Instruct 7B, which already performs strongly on benchmarks. Moreover, we offer a unified framework for interpreting a range of methods, and outline several avenues for enhancing reinforcement learning approaches.

Despite DeepSeekMath’s~noteworthy achievements in quantitative reasoning tasks, it remains less proficient in geometry and theorem-proving compared to closed models. For example, during preliminary assessments, the model was unable to solve problems involving triangles and ellipses, possibly revealing a bias in the training and fine-tuning data curation. Furthermore, DeepSeekMath’s~performance in few-shot settings lags behind models such as GPT-4, whose accuracy can improve with few-shot prompts, whereas DeepSeekMath~exhibits similar outcomes under both zero-shot and few-shot scenarios. Moving forward, we plan to refine our data selection processes to curate even higher quality pre-training corpora. Additionally, we aim to pursue the research directions proposed in Section \ref{sec:effec_rl} to further advance reinforcement learning strategies for large language models.

\end{document}